% ===================================================
% TABLES FINALES
% Pages imprimées 283-284 — vues 285-286
% ===================================================
\clearpage
\fancyhead[R]{}
\iffalse
\chapter*{Table des matières}
\addcontentsline{toc}{chapter}{Table des matières}

\begin{longtable}{@{}p{0.78\textwidth}r@{}}
Préface & 1 \\
Étymologie du nom d'Ouroux & 7 \\
Topographie & 9 \\
Période romaine & 12 \\
Période burgonde & 20 \\
Période féodale & 29 \\
Nagu & 35 \\
Arcis & 47 \\
La Carelle & 51 \\
Montaulieu & 61 \\
Gros Bois & 62 \\
Autres familles anciennes & 67 \\
Alloignet & 72 \\
Bourg, population, industrie, etc. & 75 \\
Curés d'Ouroux & 96 \\
Grande Révolution & 141 \\
Vicaires d'Ouroux & 167 \\
Prêtres nés à Ouroux & 176 \\
Église d'Ouroux & 181 \\
Écoles d'Ouroux & 231 \\
\quad École de filles & 233 \\
\quad École de garçons & 245 \\
\quad Notes sur la laïcisation et les secours accordés aux écoles libres & 257 \\
Période contemporaine & 261 \\
\quad Un maire chrétien & 262 \\
\quad Guerre 1914--1918 & 264 \\
\quad Catastrophes aériennes & 267 \\
\quad Guerre 1939--1945 -- bombardement & 268 \\
\quad Vieilles maisons & 270 \\
\quad Progrès social et matériel & 271 \\
Conclusion & 275 \\
Avenas \& son église & 277 \\
\end{longtable}

Voir à la page suivante la table des illustrations.

\fi
\clearpage
\fancypagestyle{illustrationstable}{%
    \fancyhf{}
    \fancyhead[C]{\textendash\ \thepage\ \textendash}
    \renewcommand{\headrulewidth}{0pt}
    \renewcommand{\footrulewidth}{0pt}
}
\thispagestyle{illustrationstable}
\addcontentsline{toc}{chapter}{Table des illustrations}

\begin{center}
{\large\bfseries Table des illustrations\par}
\end{center}
\vspace{-0.4em}
\enlargethispage{4\baselineskip}

\iffalse
\begin{multicols}{2}
\fontsize{7}{7.5}\selectfont
\setlength{\columnsep}{1em}
\noindent\textbf{N\textsuperscript{o}\hfill Page}

\begin{description}
\setlength{\itemsep}{-0.5ex}
\setlength{\parsep}{0pt}
\setlength{\topsep}{0pt}
\setlength{\partopsep}{0pt}
\item[1] clocher d'Ouroux \dotfill 3
\item[2] un oratoire \dotfill 7
\item[3] plan d'Ouroux \dotfill 10
\item[4] gladiateurs \dotfill 12
\item[5--6] Sts Pothin \& Irénée \dotfill 12
\item[7] château d'Avenas et vallée d'Ouroux \dotfill 14
\item[8] monnaies romaines \dotfill 17
\item[9] urnes funéraires romaines \dotfill 17
\item[10] masques tragiques \dotfill 19
\item[11] broche burgonde \dotfill 20
\item[12] aiguille mérovingienne \dotfill 20
\item[13] scène de labour (XII\textsuperscript{o} s.) \dotfill 28
\item[14] château du Sauzay (Avenas) \dotfill 29
\item[15] scène d'hommage \dotfill 32
\item[16] tour de Nagu et clocher \dotfill 35
\item[17] tour de Nagu \dotfill 37
\item[18] armes de Nagu \dotfill 39
\item[19] chevaliers au combat \dotfill 42
\item[20--21] armes Berthet \& Chantem. \dotfill 46
\item[22] maison d'Arcis \dotfill 47
\item[23--27] armoiries \dotfill 49
\item[28] château de la Carelle \dotfill 51
\item[29--31] armoiries \dotfill 56
\item[32] armes de Dubost \dotfill 57
\item[33] armes de Hautrier \dotfill 59
\item[34] armes de Laurencin \dotfill 60
\item[35] château de Gros-Bois \dotfill 62
\item[36] armes de Gros-bois \dotfill 64
\item[37] imposte en fer forgé \dotfill 65
\item[38] chapiteau au Thel \dotfill 66
\item[39] armes des Testenoire \dotfill 67
\item[40--41] scènes paysannes \dotfill 71
\item[42] ruines d'Alloignet \dotfill 72
\item[43] plan du bourg d'Ouroux \dotfill 74
\item[44] la cure et sa tour \dotfill 77
\item[45] grande rue d'Ouroux \dotfill 79
\item[46] la croix de l'étang \dotfill 82
\item[47] armes de Guillin \dotfill 85
\item[48--49] ouvrier et paysans \dotfill 87
\item[50] église de Beaujeu \dotfill 90
\item[51] le loup brûlé \dotfill 95
\item[52] cure d'Ouroux \dotfill 96
\item[53] le vieux St-Vincent \dotfill 97
\item[54] château de Gorze \dotfill 100
\item[55] descente de croix \dotfill 116
\item[56] ancien clocheton \dotfill 118
\item[57] église de St-Christophe \dotfill 123
\item[58] château de Verbust \dotfill 124
\item[59] prise de la Bastille \dotfill 141
\item[60] église de St-Jacques \dotfill 156
\item[61] église de St-Mamert \dotfill 164
\item[62] église d'Ouroux (clocher) \dotfill 181
\end{description}
\columnbreak
\noindent\textbf{N\textsuperscript{o}\hfill Page}
\begin{description}
\setlength{\itemsep}{-0.5ex}
\setlength{\parsep}{0pt}
\setlength{\topsep}{0pt}
\setlength{\partopsep}{0pt}
\item[63] abside église d'Ouroux \dotfill 183
\item[64] autel de St-Antoine \dotfill 187
\item[65] plan de l'église d'Ouroux \dotfill 190
\item[66] le tabernacle (maître-autel) \dotfill 202
\item[67] intérieur de l'église \dotfill 203
\item[68] porte principale \dotfill 204
\item[69] plan actuel de l'église \dotfill 204
\item[70--72] tableaux du 1er chapit. \dotfill 205
\item[73--76] tableaux 2ème chapit. \dotfill 206
\item[77--80] tableaux 3ème chapit. \dotfill 207
\item[80--83] tableaux 4ème chapitre \dotfill 208
\item[84--85] tab. 5 \& 6ème chapiteaux \dotfill 209
\item[86--88] tableaux 7ème chapit. \dotfill 210
\item[89--92] tableaux 8ème chapit. \dotfill 211
\item[93--96] tableaux 9ème chapit. \dotfill 212
\item[97--99] tableaux 10ème chapit. \dotfill 213
\item[100] intérieur de l'abside \dotfill 214
\item[101] vierge à la rose (autel S.V.) \dotfill 215
\item[102] la Cène (détail de vitrail) \dotfill 216
\item[103--104] vitraux (Anne, Antoine) \dotfill 217
\item[105] détail sculpture chaire \dotfill 218
\item[106] une station du chemin de + \dotfill 218
\item[107--110] vitraux modernes \dotfill 223 à 226
\item[111] fonts baptismaux \dotfill 227
\item[112] église d'Ouroux \dotfill 229
\item[113] maison des sœurs \dotfill 233
\item[114] école libre de filles \dotfill 242
\item[115] école communale de filles \dotfill 245
\item[116] mairie -- école communale \dotfill 249
\item[117] salle des fêtes \dotfill 261
\item[118] monuments aux morts \dotfill 266
\item[119] le ``Capricornus'' \dotfill 267
\item[120] église d'Avenas \dotfill 277
\item[121] clocher d'Avenas et croix \dotfill 279
\item[122] l'autel d'Avenas \dotfill 281
\item[123] inscription de l'autel \dotfill 282
\end{description}
\end{multicols}
\fi

% Cette table est produite à partir du registre central des illustrations.
% Les pages sont celles de la pagination imprimée de l'édition originale.
\begin{multicols}{2}
\fontsize{7}{7.5}\selectfont
\setlength{\columnsep}{1em}
\noindent\textbf{N\textsuperscript{o}\hfill Page}

\begin{description}
\setlength{\itemsep}{-0.5ex}
\setlength{\parsep}{0pt}
\setlength{\topsep}{0pt}
\setlength{\partopsep}{0pt}
\renewcommand{\ourouxillustrationmetadata}[3]{%
    \item[#1] #2 \dotfill #3}
% Description et page imprimee de chaque illustration de l'edition originale.
\ourouxillustrationmetadata{Illustrations/illus-01-clocher-ouroux.png}{clocher d'Ouroux et oratoire}{7}
\ourouxillustrationmetadata{Illustrations/illus-02-plan-ouroux.png}{plan d'Ouroux}{10}
\ourouxillustrationmetadata{Illustrations/illus-04-gladiateurs.png}{gladiateurs}{12}
\ourouxillustrationmetadata{Illustrations/illus-05-medaille-pothin.png}{medaillon de St-Pothin}{12}
\ourouxillustrationmetadata{Illustrations/illus-06-medaille-romaine.png}{medaillon de St-Irenee}{12}
\ourouxillustrationmetadata{Illustrations/illus-19-chateau-guillin-vallee.png}{chateau d'Avenas et vallee d'Ouroux}{14}
\ourouxillustrationmetadata{Illustrations/illus-07-amphore-monnaies.png}{monnaies romaines}{17}
\ourouxillustrationmetadata{Illustrations/illus-08-urnes-funeraires.png}{urnes funeraires romaines}{17}
\ourouxillustrationmetadata{Illustrations/illus-10-masques-tragiques.png}{masques tragiques}{19}
\ourouxillustrationmetadata{Illustrations/illus-11-broche-burgonde.png}{broche burgonde}{20}
\ourouxillustrationmetadata{Illustrations/illus-12-gaule-merovingienne.png}{aiguille merovingienne}{20}
\ourouxillustrationmetadata{Illustrations/illus-13-scene-labour.png}{scene de labour}{28}
\ourouxillustrationmetadata{Illustrations/illus-14-chateau-sauzay-avenas.png}{chateau du Sauzay}{29}
\ourouxillustrationmetadata{Illustrations/illus-15-scene-hommage.png}{scene d'hommage}{32}
\ourouxillustrationmetadata{Illustrations/illus-nagu-clocher-tour.png}{tour de Nagu et clocher}{35}
\ourouxillustrationmetadata{Illustrations/illus-nagu-tour.png}{tour de Nagu}{37}
\ourouxillustrationmetadata{Illustrations/illus-nagu-armoiries.png}{armes de Nagu}{39}
\ourouxillustrationmetadata{Illustrations/illus-nagu-chevaliers-combat.png}{chevaliers au combat}{42}
\ourouxillustrationmetadata{Illustrations/illus-nagu-armoiries-berthet-chantemerle.png}{armes Berthet et Chantemerle}{46}
\ourouxillustrationmetadata{Illustrations/illus-arcis-maison.png}{maison d'Arcis}{47}
\ourouxillustrationmetadata{Illustrations/illus-arcis-armoiries-panneau.png}{armoiries d'Arcis}{49}
\ourouxillustrationmetadata{Illustrations/illus-carelle-chateau.png}{chateau de la Carelle}{51}
\ourouxillustrationmetadata{Illustrations/illus-carelle-carrige.png}{armoiries de la Carelle}{56}
\ourouxillustrationmetadata{Illustrations/illus-carelle-roche.png}{armoiries de la Carelle}{56}
\ourouxillustrationmetadata{Illustrations/illus-carelle-magnin.png}{armoiries de la Carelle}{56}
\ourouxillustrationmetadata{Illustrations/illus-montaulieu-dubost.png}{armes de Dubost}{57}
\ourouxillustrationmetadata{Illustrations/illus-montaulieu-fautrieres.png}{armes de Hautrier}{59}
\ourouxillustrationmetadata{Illustrations/illus-montaulieu-laurencin.png}{armes de Laurencin}{60}
\ourouxillustrationmetadata{Illustrations/illus-grosbois-chateau.png}{chateau de Gros-Bois}{62}
\ourouxillustrationmetadata{Illustrations/illus-grosbois-armoiries.png}{armes de Gros-Bois}{64}
\ourouxillustrationmetadata{Illustrations/illus-grosbois-imposte.png}{imposte en fer forge}{65}
\ourouxillustrationmetadata{Illustrations/illus-grosbois-chapiteau.png}{chapiteau au Thel}{66}
\ourouxillustrationmetadata{Illustrations/illus-autres-testenoire.png}{armes des Testenoire}{67}
\ourouxillustrationmetadata{Illustrations/illus-autres-scene-ouvrier.png}{scene paysanne}{71}
\ourouxillustrationmetadata{Illustrations/illus-autres-scene-boulanger.png}{scene paysanne}{71}
\ourouxillustrationmetadata{Illustrations/illus-alloignet-ruines.png}{ruines d'Alloignet}{72}
\ourouxillustrationmetadata{Illustrations/illus-alloignet-plan-bourg.png}{plan du bourg d'Ouroux}{74}
\ourouxillustrationmetadata{Illustrations/illus-bourg-cure-tour.png}{la cure et sa tour}{77}
\ourouxillustrationmetadata{Illustrations/illus-bourg-grande-rue.png}{grande rue d'Ouroux}{79}
\ourouxillustrationmetadata{Illustrations/illus-bourg-croix-etang.png}{la croix de l'etang}{82}
\ourouxillustrationmetadata{Illustrations/illus-bourg-armoiries-guillin.png}{armes de Guillin}{85}
\ourouxillustrationmetadata{Illustrations/illus-bourg-ouvrier-paysan.png}{ouvrier et paysan}{87}
\ourouxillustrationmetadata{Illustrations/illus-bourg-eglise-beaujeu.png}{eglise de Beaujeu}{90}
\ourouxillustrationmetadata{Illustrations/illus-bourg-loup-renard.png}{le loup brule}{95}
\ourouxillustrationmetadata{Illustrations/illus-cures-vue96.png}{cure d'Ouroux}{96}
\ourouxillustrationmetadata{Illustrations/illus-cures-saint-vincent.png}{le vieux St-Vincent}{97}
\ourouxillustrationmetadata{Illustrations/illus-cures-gorze.png}{chateau de Gorze}{100}
\ourouxillustrationmetadata{Illustrations/illus-cures-descente-croix.png}{descente de croix}{116}
\ourouxillustrationmetadata{Illustrations/illus-cures-clocheton.png}{ancien clocheton}{118}
\ourouxillustrationmetadata{Illustrations/illus-cures-saint-christophe.png}{eglise de St-Christophe}{123}
\ourouxillustrationmetadata{Illustrations/illus-cures-verbust.png}{chateau de Verbust}{124}
\ourouxillustrationmetadata{Illustrations/illus-revolution-bastille.png}{prise de la Bastille}{141}
\ourouxillustrationmetadata{Illustrations/illus-revolution-jacques-arrets.png}{eglise de St-Jacques}{156}
\ourouxillustrationmetadata{Illustrations/illus-revolution-saint-mamert.png}{eglise de St-Mamert}{164}
\ourouxillustrationmetadata{Illustrations/illus-eglise-saint-antoine.jpg}{eglise d'Ouroux, clocher}{181}
\ourouxillustrationmetadata{Illustrations/illus-eglise-abside.jpg}{abside de l'eglise d'Ouroux}{183}
\ourouxillustrationmetadata{Illustrations/illus-autel-saint-antoine.jpg}{autel de St-Antoine}{187}
\ourouxillustrationmetadata{Illustrations/illus-plan-eglise.jpg}{plan de l'eglise d'Ouroux}{190}
\ourouxillustrationmetadata{Illustrations/illus-tabernacle.jpg}{tabernacle du maitre-autel}{202}
\ourouxillustrationmetadata{Illustrations/illus-eglise-interieur.jpg}{interieur de l'eglise}{203}
\ourouxillustrationmetadata{Illustrations/illus-eglise-exterieur.jpg}{porte principale}{204}
\ourouxillustrationmetadata{Illustrations/illus-plan-eglise-actuel.jpg}{plan actuel de l'eglise}{204}
\ourouxillustrationmetadata{Illustrations/illus-chapiteau-01.jpg}{tableau du chapiteau 01}{205}
\ourouxillustrationmetadata{Illustrations/illus-chapiteau-02.jpg}{tableau du chapiteau 02}{205}
\ourouxillustrationmetadata{Illustrations/illus-chapiteau-03.jpg}{tableau du chapiteau 03}{205}
\ourouxillustrationmetadata{Illustrations/illus-chapiteau-04.jpg}{tableau du chapiteau 04}{206}
\ourouxillustrationmetadata{Illustrations/illus-chapiteau-05.jpg}{tableau du chapiteau 05}{206}
\ourouxillustrationmetadata{Illustrations/illus-chapiteau-06.jpg}{tableau du chapiteau 06}{206}
\ourouxillustrationmetadata{Illustrations/illus-chapiteau-07.jpg}{tableau du chapiteau 07}{206}
\ourouxillustrationmetadata{Illustrations/illus-chapiteau-decor-06.jpg}{tableau du chapiteau 08}{207}
\ourouxillustrationmetadata{Illustrations/illus-chapiteau-08.jpg}{tableau du chapiteau 09}{207}
\ourouxillustrationmetadata{Illustrations/illus-chapiteau-09.jpg}{tableau du chapiteau 10}{207}
\ourouxillustrationmetadata{Illustrations/illus-chapiteau-10.jpg}{tableau du chapiteau 11}{207}
\ourouxillustrationmetadata{Illustrations/illus-chapiteau-11.jpg}{tableau du chapiteau 12}{208}
\ourouxillustrationmetadata{Illustrations/illus-chapiteau-12.jpg}{tableau du chapiteau 13}{208}
\ourouxillustrationmetadata{Illustrations/illus-chapiteau-13.jpg}{tableau du chapiteau 14}{208}
\ourouxillustrationmetadata{Illustrations/illus-chapiteau-14.jpg}{tableau du chapiteau 15}{208}
\ourouxillustrationmetadata{Illustrations/illus-chapiteau-15.jpg}{tableau du chapiteau 16}{209}
\ourouxillustrationmetadata{Illustrations/illus-chapiteau-16.jpg}{tableau du chapiteau 17}{209}
\ourouxillustrationmetadata{Illustrations/illus-chapiteau-17.jpg}{tableau du chapiteau 18}{210}
\ourouxillustrationmetadata{Illustrations/illus-chapiteau-18.jpg}{tableau du chapiteau 19}{210}
\ourouxillustrationmetadata{Illustrations/illus-chapiteau-19.jpg}{tableau du chapiteau 20}{210}
\ourouxillustrationmetadata{Illustrations/illus-chapiteau-20.jpg}{tableau du chapiteau 21}{211}
\ourouxillustrationmetadata{Illustrations/illus-chapiteau-21.jpg}{tableau du chapiteau 22}{211}
\ourouxillustrationmetadata{Illustrations/illus-chapiteau-22.jpg}{tableau du chapiteau 23}{211}
\ourouxillustrationmetadata{Illustrations/illus-chapiteau-23.jpg}{tableau du chapiteau 24}{211}
\ourouxillustrationmetadata{Illustrations/illus-chapiteau-24.jpg}{tableau du chapiteau 25}{212}
\ourouxillustrationmetadata{Illustrations/illus-chapiteau-25.jpg}{tableau du chapiteau 26}{212}
\ourouxillustrationmetadata{Illustrations/illus-chapiteau-26.jpg}{tableau du chapiteau 27}{212}
\ourouxillustrationmetadata{Illustrations/illus-chapiteau-27.jpg}{tableau du chapiteau 28}{212}
\ourouxillustrationmetadata{Illustrations/illus-chapiteau-28.jpg}{tableau du chapiteau 29}{213}
\ourouxillustrationmetadata{Illustrations/illus-chapiteau-29.jpg}{tableau du chapiteau 30}{213}
\ourouxillustrationmetadata{Illustrations/illus-chapiteau-30.jpg}{tableau du chapiteau 31}{213}
\ourouxillustrationmetadata{Illustrations/illus-chapiteau-31.jpg}{tableau du chapiteau 32}{214}
\ourouxillustrationmetadata{Illustrations/illus-chapiteau-decor-31.jpg}{interieur de l'abside}{214}
\ourouxillustrationmetadata{Illustrations/illus-autel-vierge-rose.jpg}{vierge a la rose}{215}
\ourouxillustrationmetadata{Illustrations/illus-sainte-anne.jpg}{vitrail Ste-Anne}{217}
\ourouxillustrationmetadata{Illustrations/illus-saint-antoine-barthelemy.jpg}{vitrail St-Antoine et St-Barthelemy}{217}
\ourouxillustrationmetadata{Illustrations/illus-chaire-bois.jpg}{detail de sculpture de la chaire}{218}
\ourouxillustrationmetadata{Illustrations/illus-chaire-decor.jpg}{station du chemin de croix}{218}
\ourouxillustrationmetadata{Illustrations/illus-vitrail-saint-martin.jpg}{vitrail de Saint-Martin}{223}
\ourouxillustrationmetadata{Illustrations/illus-vitrail-saint-louis.jpg}{vitrail de Saint-Louis}{224}
\ourouxillustrationmetadata{Illustrations/illus-vitrail-jeanne-arc.jpg}{vitrail de Sainte-Jeanne d'Arc}{225}
\ourouxillustrationmetadata{Illustrations/illus-vitrail-bernadette.jpg}{vitrail de Sainte-Bernadette}{226}
\ourouxillustrationmetadata{Illustrations/illus-fonts-baptismaux.jpg}{fonts baptismaux}{227}
\ourouxillustrationmetadata{Illustrations/illus-eglise-final.jpg}{eglise d'Ouroux}{229}
\ourouxillustrationmetadata{Illustrations/illus-ecole-maison-soeurs.jpg}{maison des soeurs}{233}
\ourouxillustrationmetadata{Illustrations/illus-ecole-filles-berloty.jpg}{ecole libre de filles}{242}
\ourouxillustrationmetadata{Illustrations/illus-ecole-filles-1914.jpg}{ecole communale de filles}{245}
\ourouxillustrationmetadata{Illustrations/illus-salle-fetes.jpg}{salle des fetes}{261}
\ourouxillustrationmetadata{Illustrations/illus-plaque-soldats.jpg}{plaque des soldats}{266}
\ourouxillustrationmetadata{Illustrations/illus-monument-morts.jpg}{monument aux morts}{266}
\ourouxillustrationmetadata{Illustrations/illus-catastrophe-aerienne.jpg}{le Capricornus}{267}
\ourouxillustrationmetadata{Illustrations/illus-maisons-anciennes.jpg}{maisons anciennes}{270}
\ourouxillustrationmetadata{Illustrations/illus-avenas-eglise.jpg}{eglise d'Avenas}{277}
\ourouxillustrationmetadata{Illustrations/illus-avenas-croix.jpg}{clocher d'Avenas et croix}{279}
\ourouxillustrationmetadata{Illustrations/illus-autel-avenas-bas-relief.jpg}{autel d'Avenas}{281}
\ourouxillustrationmetadata{Illustrations/illus-autel-avenas-inscription.jpg}{inscription de l'autel}{282}

\end{description}
\end{multicols}

\begin{center}
\footnotesize
\textit{Enluminures diverses... passim.}

Ce livre a été tiré à 200 exemplaires, tous numérotés en seconde page.
\end{center}
